\documentclass[a4paper]{article}

% Español (México) con XeLaTeX: fontspec para fuentes Unicode, polyglossia
% para la división silábica y los nombres del documento en la variante mexicana.
\usepackage{fontspec}
\usepackage{polyglossia}
\setdefaultlanguage[variant=mexican]{spanish}
% Los nombres chinos van como «pinyin (original)»: xeCJK da a los caracteres
% Han su propia fuente; sin ella XeLaTeX los omite con un aviso.
% FandolSong no cubre algunos caracteres de nombres (U+73FA); WenQuanYi Zen Hei sí.
\usepackage[AutoFallBack=true]{xeCJK}
\setCJKmainfont{FandolSong-Regular.otf}
\setCJKfallbackfamilyfont{\CJKrmdefault}{WenQuanYi Zen Hei}
% polyglossia no traduce los nombres de listings.
\AtBeginDocument{\providecommand{\lstlistingname}{}\renewcommand{\lstlistingname}{Listado}%
  \providecommand{\lstlistlistingname}{}\renewcommand{\lstlistlistingname}{Índice de listados}}
\usepackage{amsmath,amssymb}
\usepackage{graphicx}
\usepackage[margin=2.5cm]{geometry}
\usepackage[most]{tcolorbox}
\usepackage{etoolbox}
\usepackage{subcaption}
\usepackage{listings}
\usepackage{hyperref}
\usepackage{booktabs}
\usepackage{float}

\newtcolorbox{knowledgebox}[1]{
    enhanced, colback=blue!5!white, colframe=blue!75!black, colbacktitle=blue!75!black,
    coltitle=white, fonttitle=\bfseries, title={#1}, boxrule=1pt, sharp corners
}

\newtcolorbox{importantbox}[1]{
    enhanced, colback=yellow!10!white, colframe=yellow!80!black, colbacktitle=yellow!80!black,
    coltitle=black, fonttitle=\bfseries, title={#1}, sharp corners
}

\newtcolorbox{warningbox}[1]{
    enhanced, colback=red!5!white, colframe=red!75!black, colbacktitle=red!75!black,
    coltitle=white, fonttitle=\bfseries, title={#1}, sharp corners
}

\lstset{
    basicstyle=\ttfamily\small,
    keywordstyle=\color{blue},
    stringstyle=\color{red!60!black},
    commentstyle=\color{green!60!black},
    breaklines=true,
    frame=single,
    numbers=left,
    numberstyle=\tiny\color{gray},
    captionpos=b,
    extendedchars=true
}

\newcommand{\notetitle}{CS224R Lecture 19: repaso y síntesis de Q-Learning}
\newcommand{\noteauthors}{Elaborado por Anakiette y el equipo docente de Stanford CS224R}
\newcommand{\notedate}{3 de abril de 2026}
\newcommand{\videochannel}{Stanford Online}
\newcommand{\videopublishdate}{2025}
\newcommand{\videoduration}{aproximadamente 60 minutos}
\newcommand{\videourl}{https://cs224r.stanford.edu/}
\newcommand{\videocoverpath}{cover.jpg}
\newcommand{\repourl}{https://github.com/hqhq1025/ai-course-notes}


% Footer with repo info
\usepackage{fancyhdr}
\pagestyle{fancy}
\fancyhf{}
\fancyfoot[C]{\thepage}
\fancyfoot[R]{\footnotesize\color{gray}\href{\repourl}{AI Course Notes}}
\renewcommand{\headrulewidth}{0pt}
\renewcommand{\footrulewidth}{0pt}

\begin{document}
\begin{titlepage}
\centering
{\Large Notas del curso\par}
\vspace{1.2cm}
{\huge\bfseries \notetitle\par}
\vspace{0.8cm}
{\large \noteauthors\par}
\vspace{0.3cm}
{\large \notedate\par}
\vspace{1.1cm}

\ifdefempty{\videocoverpath}{
{\small No se proporcionó imagen de portada.\par}
}{
\includegraphics[width=0.82\textwidth,height=0.45\textheight,keepaspectratio]{\videocoverpath}\par
}

\vfill
\begin{tcolorbox}[width=0.9\textwidth, colback=black!2!white, colframe=black!60, sharp corners]
\textbf{Autor o canal del video}: \videochannel\par
\textbf{Fecha de publicación}: \videopublishdate\par
\textbf{Duración del video}: \videoduration\par
\textbf{Ponente}: Anakiette (TA)\par
\textbf{Página del curso}: \href{https://cs224r.stanford.edu/}{cs224r.stanford.edu}\par
\textbf{Página del proyecto}: \href{\repourl}{\nolinkurl{\repourl}}\par
\end{tcolorbox}
\end{titlepage}

\tableofcontents
\newpage

%% --- Inicio del contenido --- %%

\section{Repaso de los fundamentos de los MDP}

Esta sesión la imparte el asistente de docencia Anakiette como repaso general del curso: parte de los fundamentos de los MDP y recorre la cadena teórica completa de Q-learning.

\textit{``We're going to start off by just doing a very very brief overview of MDPs.''} Esta frase inicial es a la vez una indicación y una guía para el ritmo de estudio: recuerda que, dentro de la cadena teórica completa, primero se establece el sistema de notación formal y después se profundiza en los detalles nivel por nivel.

\subsection{Procesos de decisión de Markov}

\begin{importantbox}{La quíntupla de un MDP}
Un MDP se define mediante $(\mathcal{S}, \mathcal{A}, P, R, \gamma)$:
\begin{itemize}
    \item $\mathcal{S}$: espacio de estados
    \item $\mathcal{A}$: espacio de acciones
    \item $P(s' | s, a)$: probabilidad de transición de estados
    \item $R(s, a)$ o $R(s, a, s')$: función de recompensa
    \item $\gamma \in [0, 1)$: factor de descuento
\end{itemize}
\textbf{Propiedad de Markov}: el siguiente estado depende solo del estado y la acción actuales, no del historial.
\end{importantbox}

\subsection{Función de valor y función Q}

\begin{importantbox}{Definiciones fundamentales}
\textbf{Función de valor de estado}: retorno acumulado esperado al seguir la política $\pi$ desde el estado $s$
$$V^\pi(s) = \mathbb{E}_\pi\left[\sum_{t=0}^{\infty} \gamma^t R(s_t, a_t) \mid s_0 = s\right]$$

\textbf{Función de valor de acción} (función Q): retorno acumulado esperado al ejecutar la acción $a$ en el estado $s$ y seguir después la política $\pi$
$$Q^\pi(s, a) = \mathbb{E}_\pi\left[\sum_{t=0}^{\infty} \gamma^t R(s_t, a_t) \mid s_0 = s, a_0 = a\right]$$

Relación entre ambas: $V^\pi(s) = \sum_a \pi(a|s) Q^\pi(s, a)$
\end{importantbox}

\subsection{Ecuaciones de Bellman}

\begin{importantbox}{Ecuación de esperanza de Bellman}
$$Q^\pi(s, a) = R(s, a) + \gamma \sum_{s'} P(s'|s,a) \sum_{a'} \pi(a'|s') Q^\pi(s', a')$$

\textbf{Ecuación de optimalidad de Bellman}:
$$Q^*(s, a) = R(s, a) + \gamma \sum_{s'} P(s'|s,a) \max_{a'} Q^*(s', a')$$
\end{importantbox}

\subsection{Iteración de políticas e iteración de valores}

\begin{knowledgebox}{Dos enfoques iterativos complementarios}
La iteración de políticas (policy iteration) y la iteración de valores (value iteration) pueden verse como dos formas de llevar a la práctica la ecuación de optimalidad de Bellman: la primera optimiza de manera alternada la política y la función de valor; la segunda itera directamente sobre el valor óptimo. En cuanto a la convergencia, ambas dependen de que la iteración de Bellman sea una contracción; sin embargo, la primera actualiza la política con mayor precisión en cada ronda, mientras que la segunda requiere un cálculo más sencillo en cada paso.
\end{knowledgebox}

\begin{table}[H]
\centering
\begin{tabular}{ll}
\toprule
Método & Características \\
\midrule
Iteración de políticas & Construye explícitamente la política de comportamiento (alineada con la política anterior) y cambia la política solo cuando la función de valor converge \\
Iteración de valores & Itera directamente $V(s) \leftarrow \max_a [R+\gamma V(s')]$, sin almacenar una política explícita \\
Relación entre ambas & La iteración de valores puede verse como un caso particular de la iteración de políticas (la política se extrae solo al final) \\
\bottomrule
\end{tabular}
\caption{Comparación de los operadores iterativos más comunes en MDP discretos}
\label{tab:mdp-iteration}
\end{table}

\subsection{Resumen de la sección}

Los MDP proporcionan el marco formal del RL. La función de valor y la función Q son las herramientas fundamentales para evaluar la calidad de una política, y las ecuaciones de Bellman describen la relación recursiva entre ellas.
\section{Q-Learning tabular}

\subsection{El algoritmo Q-Learning}

\begin{importantbox}{Regla de actualización de Q-Learning}
Q-Learning es un algoritmo off-policy y model-free que aprende directamente la función Q óptima:
$$Q(s, a) \leftarrow Q(s, a) + \alpha \left[r + \gamma \max_{a'} Q(s', a') - Q(s, a)\right]$$
donde:
\begin{itemize}
    \item $\alpha$: tasa de aprendizaje
    \item $r + \gamma \max_{a'} Q(s', a')$: TD target (objetivo de diferencia temporal)
    \item $r + \gamma \max_{a'} Q(s', a') - Q(s, a)$: TD error
\end{itemize}
\end{importantbox}

\begin{knowledgebox}{El carácter off-policy de Q-Learning}
Q-Learning es off-policy porque usa $\max_{a'}$ para calcular el TD target, sin importar cuál sea la política de comportamiento (la política con la que se recopilan los datos). Esto significa que podemos explorar con cualquier política (por ejemplo, $\epsilon$-greedy) y, al mismo tiempo, aprender la política óptima.
\end{knowledgebox}

\subsection{Estrategias de exploración y escalamiento de recompensas}

\begin{knowledgebox}{Estrategias de exploración comunes}
\begin{itemize}
    \item $\epsilon$-greedy: la mayor parte del tiempo elige la acción óptima y explora con probabilidad $1-\epsilon$; es adecuada para espacios de acciones discretos.
    \item Boltzmann/softmax: controla la aleatoriedad mediante una temperatura, lo que facilita tener en cuenta el valor relativo de las distintas acciones.
    \item Upper Confidence Bound (UCB): incorpora intervalos de confianza y aumenta automáticamente el peso de la exploración en los estados con pocas muestras.
\end{itemize}
\end{knowledgebox}

\begin{importantbox}{Escalamiento de recompensas y recompensas desbalanceadas}
Q-Learning es sensible a la distribución de las recompensas: una escala demasiado grande hace que el TD error se dispare, y una demasiado pequeña hace que el gradiente se desvanezca. Las prácticas habituales incluyen normalizar las recompensas, restar la reward promedio como línea base y aplicar ingeniería de forma a las recompensas dispersas (dense reward shaping). En la práctica, la instrumentation correspondiente debe registrar el TD error y la distribución de reward de cada par estado-acción, para ajustar a tiempo la estrategia de escalamiento.
\end{importantbox}

\begin{warningbox}{Convergencia engañosa con recompensas dispersas}
Cuando muchos estados nunca reciben una recompensa positiva, el agent puede converger a una política óptima local dentro de la región de recompensas negativas. Se recomienda revisar periódicamente la proporción de muestras positivas y negativas en el replay buffer, y aumentar la densidad mediante hindsight experience replay o recompensas auxiliares multitarea.
\end{warningbox}

\subsection{Condiciones de convergencia}

\begin{importantbox}{Garantía de convergencia de Q-Learning}
En el caso tabular (espacios de estados y de acciones finitos), Q-Learning tiene garantizada la convergencia a $Q^*$ si se cumplen las siguientes condiciones:
\begin{enumerate}
    \item Cada par estado--acción se visita un número infinito de veces
    \item La tasa de aprendizaje cumple las condiciones de Robbins-Monro: $\sum_t \alpha_t = \infty$, $\sum_t \alpha_t^2 < \infty$
\end{enumerate}
\end{importantbox}

\subsection{Resumen de la sección}

Q-Learning tabular aproxima directamente la función Q óptima mediante actualizaciones TD y cuenta con garantías teóricas de convergencia.
\section{Fitted Q-Iteration y redes Q profundas}

\subsection{De las tablas a la aproximación de funciones}

Cuando el espacio de estados es continuo o muy grande, no es posible mantener una tabla de valores Q. La solución consiste en parametrizar la función Q con un aproximador de funciones, por ejemplo, una red neuronal.

\begin{importantbox}{Fitted Q-Iteration (FQI)}
FQI convierte Q-Learning en un problema de aprendizaje supervisado:
\begin{enumerate}
    \item Con la red Q actual se calcula el TD target: $y_i = r_i + \gamma \max_{a'} Q_{\theta^-}(s_i', a')$
    \item Se toma $(s_i, a_i)$ como entrada y $y_i$ como etiqueta, y se entrena la red Q:
    $$\min_\theta \sum_i \left(Q_\theta(s_i, a_i) - y_i\right)^2$$
    \item Se actualiza la red objetivo $\theta^- \leftarrow \theta$
    \item Se repite
\end{enumerate}
\end{importantbox}

\subsection{Técnicas fundamentales de DQN}

\begin{importantbox}{Innovaciones de DQN (Deep Q-Network)}
DQN introdujo dos técnicas fundamentales para que el Q-Learning profundo funcione de manera estable:
\begin{enumerate}
    \item \textbf{Experience Replay}: la experiencia histórica se almacena en un replay buffer y, durante el entrenamiento, se muestrean mini-batches al azar. Esto rompe la correlación temporal de los datos y mejora el aprovechamiento de las muestras.
    \item \textbf{Target Network}: se usa una red objetivo $Q_{\theta^-}$, actualizada de forma periódica, para calcular el TD target. Esto reduce la inestabilidad que provoca el problema del «objetivo móvil».
\end{enumerate}
\end{importantbox}

\begin{warningbox}{Inestabilidad de Q-Learning con aproximación de funciones}
A diferencia del caso tabular, Q-Learning con aproximación de funciones \textbf{no garantiza la convergencia}. La propiedad de contracción de la actualización de Bellman puede dejar de cumplirse con aproximación de funciones, lo que provoca:
\begin{itemize}
    \item Divergencia de los valores Q (acumulación del overestimation bias)
    \item Entrenamiento inestable (oscillation o divergence)
    \item Sensibilidad a los hiperparámetros
\end{itemize}
Por eso se necesitan técnicas de estabilización como la target network y la experience replay.
\end{warningbox}

\subsection{Mejoras de DQN}

\begin{itemize}
    \item \textbf{Double DQN}: la red actual selecciona la acción y la red objetivo evalúa su valor, lo que reduce la overestimation
    \item \textbf{Dueling DQN}: descompone la función Q en $V(s) + A(s, a)$ y estima por separado el valor del estado y la ventaja
    \item \textbf{Prioritized Experience Replay}: muestrea con prioridad según la magnitud del TD error, para concentrarse en la experiencia de la que más se puede aprender
\end{itemize}

\subsection{Resumen de la sección}

El paso central del Q-Learning tabular a DQN es la aproximación de funciones. La experience replay y la target network son las técnicas fundamentales para que el entrenamiento sea estable.

\section{De DQN a Actor-Critic}

\subsection{Limitaciones de Q-Learning}

\begin{warningbox}{Dificultades de Q-Learning en espacios de acciones continuos}
Q-Learning necesita calcular $\max_{a'} Q(s', a')$. En un espacio de acciones discreto se pueden recorrer todas las acciones. Pero en un espacio de acciones continuo, esta maximización se convierte en un problema de optimización continua, con un costo computacional muy alto.
\end{warningbox}

\subsection{Métodos Actor-Critic}

Solución: aprender al mismo tiempo la política (actor) y la función de valor (critic):
\begin{itemize}
    \item \textbf{Actor}: $\pi_\theta(a|s)$, política parametrizada que produce las acciones.
    \item \textbf{Critic}: $Q_\phi(s, a)$, evalúa qué tan buenas son las acciones que elige el actor.
    \item El actor se actualiza con gradiente de política y el critic, con aprendizaje TD.
\end{itemize}

\subsection{Resumen de la sección}

Los métodos Actor-Critic combinan las ventajas de los métodos de gradiente de política y de los métodos basados en funciones de valor, y son el paradigma dominante para tratar espacios de acciones continuos.

\section{Q-Learning: práctica de sistemas y monitoreo}

\subsection{Métricas de monitoreo y calidad de los datos}

\begin{knowledgebox}{Métricas de observación clave}
\begin{itemize}
    \item Cobertura del replay buffer: mide si las últimas $N$ transiciones cubren los pares estado-acción frecuentes.
    \item Distribución del TD error: monitorear la desviación estándar y el valor máximo del TD error permite detectar con anticipación una explosión numérica.
    \item Action distribution drift: verificar si la frecuencia de muestreo de la política de comportamiento se concentra en unas pocas acciones, para evitar que el entrenamiento derive hacia una política estrecha.
\end{itemize}
\end{knowledgebox}

Para lograr un despliegue estable, el equipo de ingeniería debe exponer estas métricas en el tablero de entrenamiento y definir umbrales de alerta con base en datos de referencia. La experiencia que se incorpora debe etiquetarse (por ejemplo, desafiante, reward escasa, cambio de entorno) para facilitar el análisis posterior.

\begin{importantbox}{Advertencia sobre la cobertura y la distribución del Replay Buffer}
Si el buffer se llena justamente con experiencias exitosas tempranas, las muestras de la exploración posterior no pueden entrar al entrenamiento, lo que provoca que el agent converja prematuramente. Se recomienda configurar \texttt{min\_size} y \texttt{max\_size}, revisar periódicamente las marcas de tiempo mediante la interfaz de muestreo y, cuando sea necesario, usar \texttt{prioritized replay} o \texttt{rank-based} para volver a incentivar el muestreo.
\end{importantbox}

\subsection{Estrategias de repetición de experiencia y rebalanceo de datos}

\begin{knowledgebox}{Técnicas de rebalanceo de Experience Replay}
\begin{itemize}
    \item Prioritized Replay: usa $|\delta|^\alpha$ como peso de muestreo, donde $\delta$ es el TD error y $\alpha$ controla la intensidad de la preferencia.
    \item Importance Sampling Correction: aplica una compensación proporcional a las muestras priorizadas para evitar el policy bias.
    \item Diversity Buffer: conserva fragmentos de experiencia provenientes de distintas subtareas o de distintas seed de simulación, para evitar que el entrenamiento quede atrapado en una narrow distribution.
\end{itemize}
\end{knowledgebox}

No basta con tratar el replay buffer como una caja negra; se recomienda registrar, después de cada iteración de entrenamiento, la proporción de muestras nuevas y antiguas que ingresan y su reward promedio. Si la proporción de muestras nuevas es menor que 30\%, la eficiencia de la exploración está disminuyendo y es necesario aumentar la tasa de aprendizaje o epsilon.

\begin{warningbox}{Dependencia excesiva del Replay buffer en datos antiguos}
Conforme avanza el entrenamiento, los datos antiguos pueden quedar obsoletos, sobre todo en entornos no estacionarios. Al continuar el entrenamiento hay que vigilar el timestep promedio del buffer; si más del 50\% de los datos proviene de los últimos diez minutos, se puede activar una limpieza cíclica o restablecer la prioridad a un valor alto, para evitar que la política antigua domine el gradiente.
\end{warningbox}

\subsection{Resumen de la sección}

Esta sección enfatiza la observabilidad en la construcción de sistemas de Q-Learning y la salud del replay buffer: todo reentrenamiento a gran escala debe establecer guardrail en tres dimensiones: la recolección de datos, la distribución de muestreo y el TD error.
\section{Variantes de la red Q profunda y comparación de entrenamiento}

\subsection{Análisis de las variantes principales}

\begin{table}[H]
\centering
\begin{tabular}{lll}
\toprule
Variante & Mecanismo central & Efecto \\
\midrule
Double DQN & La red actual elige la acción y la red objetivo la evalúa & Reduce la overestimation y mejora la estabilidad \\
Dueling DQN & Descompone $Q=V+A$ & Aprende más rápido el valor del estado; la ventaja de las acciones puede compartir parámetros \\
Prioritized Replay & Muestreo ponderado según el TD error & Recupera pronto las transiciones poco frecuentes pero decisivas \\
Rainbow & Combina varias técnicas de mejora (NoisyNet, Multi-step, Distributional) & Mayor cobertura, pero el ajuste de hiperparámetros es más complejo \\
\bottomrule
\end{tabular}
\caption{Variantes de DQN y su efecto en la ingeniería}
\end{table}

\begin{knowledgebox}{Escenarios adecuados para la combinación de varias técnicas}
No todos los escenarios necesitan el conjunto completo de mejoras de Rainbow. En entornos de baja capacidad, Double DQN combinado con Dueling basta para contener la sobreestimación. En tareas de Atari con abundantes recursos de cómputo o en la conducción simulada, los retornos de varios pasos de Rainbow y NoisyNet aportan una mayor capacidad de generalización.
\end{knowledgebox}

\subsection{Comparación de costos de entrenamiento e inferencia}

\begin{importantbox}{Modelo de costos: entrenamiento frente a inferencia}
El entrenamiento se concentra en el muestreo y la optimización: $Cost_{train} \approx N_{\text{steps}} \times (^{\text{forward}}C + ^{\text{backward}}C)$; la inferencia, en cambio, busca un equilibrio entre rendimiento y latencia: $Cost_{infer} \approx C_{\text{forward}} + C_{\text{action selection}}$. Si se lleva DQN a dispositivos embebidos, conviene considerar la cuantización o la poda para reducir el cómputo.
\end{importantbox}

\begin{knowledgebox}{Desglose de la pila de inferencia}
Un inference pipeline suele incluir: la percepción del entorno (state encoding), la pasada hacia adelante de la red Q, la selección de la acción con $\arg\max$, el filtrado de acciones (por ejemplo, con restricciones previas) y las verificaciones de seguridad en la interfaz con el simulador. En la coordinación multi-agent también hay que controlar los retrasos de comunicación y de sincronización.
\end{knowledgebox}

\begin{warningbox}{Cuantización excesiva en la inferencia de baja latencia}
Para optimizar la latency de inferencia, los equipos de ingeniería suelen cuantizar o podar la red; sin embargo, si no se preserva la ruta crítica (por ejemplo, la rama de advantage de la última capa), el orden de los valores Q puede alterarse. Se recomienda mantener dos canales: la versión cuantizada para la inferencia y un respaldo en full precision para la periodic calibration.
\end{warningbox}

\subsection{Resumen de la sección}

Esta sección, mediante la comparación de las variantes de DQN y el análisis del modelo de costos, señala cómo elegir las técnicas de mejora y la combinación de operadores adecuadas según las distintas restricciones de recursos.
\section{Casos de aplicación y evolución cronológica}

\subsection{Contexto de los benchmarks y línea de tiempo}

\begin{knowledgebox}{Línea de tiempo de la evolución de Q-Learning}
\begin{itemize}
    \item 1992: Watkins demuestra la convergencia teórica de Q-Learning.
    \item 2013: Mnih et al. logran control end-to-end en Atari con DQN.
    \item 2015-2017: Double DQN, Dueling, Prioritized Replay, Rainbow y otras variantes van perfeccionando el método.
    \item Después de 2016: DQN se extiende a tareas de control continuo como MuJoCo y Roboschool, y empiezan a popularizarse el RL multi-agent y el entrenamiento adversarial.
\end{itemize}
\end{knowledgebox}

Al proponer un benchmark siempre hay que elegir entre usar un entorno sencillo para obtener reproducibilidad o un entorno complejo para aproximarse a las aplicaciones reales. El repaso del curso de Stanford subrayó el trade-off entre el control que ofrece el primero y la representatividad del segundo.

\subsection{Lecciones de los casos}

\begin{importantbox}{Revisión del caso de Atari}
En la tarea Atari Breakout, el Agent extrae información de movimiento mediante frame stacking, mientras que frame-skip mejora la eficiencia del entrenamiento. Al modificar el reward shaping (por ejemplo, al aumentar la recompensa por golpear un brick), la estrategia del agent cambió de forma notable, lo que muestra que ajustar las reward scales puede provocar fácilmente estrategias nuevas.
\end{importantbox}

\begin{warningbox}{Sesgo de observación en aplicaciones reales}
En las tareas de robótica, el ruido de los sensor reales y las diferencias con los sensor simulados pueden hacer que falle la transferencia sim-to-real. Aunque se alcance una puntuación alta en simulación, es indispensable aplicar domain randomization y calibrar con datos reales antes del despliegue.
\end{warningbox}

\subsection{Resumen de la sección}

Esta sección combina la evolución de los benchmarks con casos clásicos para recordarnos que, aun cuando dominemos los detalles del algoritmo, debemos cuidarnos de la brecha de información entre la simulación y el mundo real.
\section{Fallas comunes y gobernanza}

\subsection{Rasgos comunes de las fallas}

\begin{warningbox}{Falla de generalización durante el entrenamiento}
Aunque obtenga una reward alta en el entorno de entrenamiento, el agent puede fallar ante un estado inicial ligeramente modificado. La causa suele ser que el replay buffer depende en exceso de un pequeño grupo de transiciones con reward alta y no logra cubrir el espacio de estados. Esto puede mitigarse con un entropy bonus o con replay multitarea.
\end{warningbox}

\begin{knowledgebox}{Aspectos de gobernanza}
\begin{itemize}
    \item Policy drift: hay que conservar muestras de la política dev como referencia dentro de la ventana.
    \item Reward hacking: en el aprendizaje por refuerzo, la reward shaping exige cautela; no se debe modificar la reward de forma simple de modo que el agent encuentre un shortcut.
    \item Transparencia de los datos: registrar el slice de datos usado en cada actualización y cotejar periódicamente los registros de comportamiento con las áreas de negocio.
\end{itemize}
\end{knowledgebox}

\subsection{Mecanismos de intervención}

\begin{importantbox}{Safe Interruptibility e intervención humana}
Durante el entrenamiento, mantener la interruptibility significa que, cuando una persona detecta una anomalía, puede detener o reiniciar el agent de forma segura sin provocar una explosión de la recompensa. En la práctica se configura un \texttt{kill switch} y, al mismo tiempo, se registra el último TD error para ayudar a localizar la causa.
\end{importantbox}

\begin{warningbox}{Fragmentación del entrenamiento por intervención excesiva}
Interrumpir el entrenamiento con frecuencia hace que los datos del replay buffer provengan de políticas distintas, lo que rompe el supuesto off-policy. Se recomienda ejecutar un warm-start offline training después de cada intervención, para que la nueva política tenga tiempo suficiente de converger.
\end{warningbox}

\subsection{Resumen de la sección}

En el plano de la gobernanza hay que equilibrar el entrenamiento automático con la supervisión humana; los mecanismos de intervención deben centrarse en las guardrails y evitar la fragmentación por entrenamientos repetidos.
\section{Glosario y herramientas}

\subsection{Consulta rápida de términos clave}

\begin{table}[H]
\centering
\begin{tabular}{lp{10cm}}
\toprule
Término & Definición \\
\midrule
TD error & $r + \gamma \max_{a'} Q(s', a') - Q(s, a)$, sirve para medir la magnitud de la actualización de la estimación actual de Q \\
Prioritized Experience Replay & Muestrea según el TD error e introduce la IS correction \\
Target Network & Red que se sincroniza periódicamente y estabiliza el TD target \\
State Action Coverage & Distribución de los distintos pares state-action en el replay buffer \\
\bottomrule
\end{tabular}
\caption{Términos principales de esta clase}
\end{table}

\begin{knowledgebox}{Inventario de la cadena de herramientas}
Entre las herramientas de uso común están: OpenAI Baselines (plantilla básica de DQN), RLlib (admite múltiples agents), MLFlow + TensorBoard (registro de métricas), WeTest o Supersonic (validación de la cuantización en el edge).
\end{knowledgebox}

\subsection{Resumen de la sección}

Dominar a la vez el glosario y la cadena de herramientas ayuda a pasar de la teoría a la práctica de ingeniería; en particular, evita repetir la exploración por ensayo y error al reproducir un benchmark y al poner un modelo en producción.

 \subsection{Resumen de la sección}

Dominar a la vez el glosario y la cadena de herramientas ayuda a pasar de la teoría a la práctica de ingeniería; en particular, evita repetir la exploración por ensayo y error al reproducir un benchmark y al poner un modelo en producción.

\section{Evaluación de políticas y diseño experimental}

\subsection{Métricas fuera de línea y comparación}

La evaluación fuera de línea suele realizarse en un entorno holdout distinto del conjunto de entrenamiento, para asegurar que la política no se haya sobreajustado. Las métricas más comunes incluyen la reward acumulada, la longitud del episode, la estabilidad (desviación estándar) y la sample efficiency. La siguiente tabla muestra qué observa cada una de estas métricas:

\begin{table}[H]
\centering
\begin{tabular}{lp{8cm}}
\toprule
Métrica & Observaciones \\
\midrule
Reward acumulada & Refleja el grado en que la política completa la tarea a partir de un punto de inicio fijo \\
Episode length & Evalúa si aprendió una ruta más concisa / eficiente \\
Sample efficiency & Magnitud de la mejora en reward por unidad de interacción con el entorno \\
Estabilidad & Fluctuación de la reward entre distintas seeds; si es demasiado grande, implica riesgos en el despliegue \\
\bottomrule
\end{tabular}
\caption{Métricas habituales de evaluación fuera de línea}
\end{table}

\begin{importantbox}{Equilibrio entre sample efficiency y reward ceiling}
En ocasiones, el beneficio marginal de aumentar la reward acumulada es muy bajo, pero la sample efficiency todavía no se ha estabilizado, lo que indica que el agent sigue explorando de forma repetida. Se recomienda dar seguimiento a ambas en la curva de iteraciones y, durante el periodo de plateau, introducir una exploration más intensa o una auxiliary task.
\end{importantbox}

\subsection{Comparación en línea e intervención}

En el entorno de producción, una política nueva puede probarse mediante una comparación del tipo champion/challenger o A/B. Debe garantizarse que el challenger opere solo con una fracción pequeña del tráfico y tener preparado un plan de reversión.

\begin{knowledgebox}{A/B vs Champion/ Challenger}
A/B: el tráfico en línea se dirige directamente al challenger según una proporción; es adecuado cuando la diferencia entre la behavior policy y el baseline es pequeña. Champion/Challenger: el challenger se ejecuta en un entorno shadow y, mediante offline evaluation, se decide si se le hace promote.
\end{knowledgebox}

\begin{warningbox}{Desplazamiento de la distribución en la evaluación}
En cuanto existe una diferencia en el ruido de observación entre el entorno de entrenamiento y el de despliegue, la señal de las métricas fuera de línea puede distorsionarse. En sistemas críticos para la seguridad (como la conducción autónoma), debe darse prioridad a construir una biblioteca de scenarios para realizar stress tests.
\end{warningbox}

\subsection{Resumen de la sección}

Esta sección repasó el sistema de evaluación de Q-Learning, subrayó la importancia de atender a la vez la reward acumulada, la eficiencia de muestreo y la estabilidad, y recomendó establecer pruebas comparativas y stress tests antes del despliegue.
\section{Interpretabilidad, depuración y automatización}

\subsection{Visualización y diagnóstico de los valores Q}

Durante el entrenamiento, se puede observar con un heatmap o un contour cómo se distribuyen los valores Q según cambian las acciones y los estados. Si una acción recibe siempre un puntaje Q muy alto y, aun así, nunca se elige, eso puede indicar que la implementación de la política de comportamiento (por ejemplo, $\epsilon$-greedy) está afectada por un bug.

\begin{knowledgebox}{Elementos de las herramientas de diagnóstico}
La información importante incluye: los gradientes de Q, la curva del TD error y la action frequency en el Replay buffer. Para mostrar gráficas de series de tiempo se pueden usar TensorBoard y MLFlow.
\end{knowledgebox}

\subsection{Automatización del pipeline de entrenamiento}

Se recomienda incorporar comprobaciones automáticas en el ciclo de entrenamiento: cada cierto número de epochs, hacer una offline evaluation, sincronizar las metrics con el dashboard y calcular automáticamente la sample efficiency. El pipeline de entrenamiento también debe admitir `checkpoint+resume` y `metric gating`; este último pausa el entrenamiento automáticamente cuando bajan las métricas clave.

\begin{importantbox}{Centinela del Replay buffer}
Un modelo ligero calcula en tiempo real la coverage y las reward statistics del buffer, y emite una alerta en cuanto detecta un desplazamiento brusco de la distribución. Este centinela permite evitar que, por recolectar datos one-hot desde las primeras etapas, el entrenamiento posterior colapse.
\end{importantbox}

\begin{warningbox}{Dependencia excesiva de la automatización en métricas históricas}
Si el gating solo considera el promedio histórico de la reward, puede pasar por alto el recent decay. Se recomienda introducir una sliding window version de las métricas y escalar las anomalías a un operator para que las confirme manualmente.
\end{warningbox}

\subsection{Resumen de la sección}

La interpretabilidad y la automatización, trabajando en paralelo, son la base del entrenamiento a gran escala; organizar bien las vistas de diagnóstico y los mecanismos automáticos de gating reduce de manera considerable el número de reinicios.
\section{Direcciones futuras y problemas abiertos}

\subsection{Dependencias de largo plazo y composición de tareas}

En su forma original, Q-Learning se vuelve ineficiente ante un horizon largo o una sparse reward. Los métodos modernos tienden a combinar actualizaciones Q de múltiples pasos (multi-step TD), una hierarchical policy y un planner externo.

\begin{importantbox}{El pipeline de la composición de tareas}
Dividir una tarea única en dos niveles, Macro y Micro (por ejemplo, la meta policy se encarga de elegir la subtarea y la micro policy, del control concreto), permite atenuar la dificultad de la long-term credit assignment.
\end{importantbox}

\subsection{Gobernanza y transferibilidad}

\begin{knowledgebox}{Lista de problemas abiertos}
\begin{itemize}
    \item Sim-to-real gap: cómo lograr que los datos reales del replay buffer tengan un papel más importante.
    \item Reward hacking: la reward shaping de largo plazo puede introducir atajos (shortcut), por lo que conviene revisar periódicamente la fórmula de la recompensa (reward formula).
    \item Estados multimodales: cuando el robot state incluye imágenes, LiDAR y lenguaje, la complejidad del pipeline de procesamiento de la Q network aumenta de forma abrupta.
\end{itemize}
\end{knowledgebox}

\begin{warningbox}{Gestión del drift durante el despliegue}
Aunque la política tenga un desempeño excelente en la simulación offline, el Observation drift o el action constraint drift pueden provocar fallas en producción. Se recomienda configurar una sanity check periódica (por ejemplo, evaluate on replayed expert rollouts).
\end{warningbox}

\subsection{Resumen de la sección}

Esta sección señala que el campo de Q-Learning sigue en evolución y que las dependencias de largo plazo, la composición de tareas, y la gobernanza y la auditoría son las principales direcciones del trabajo futuro.
\section{Síntesis y ampliación}

\subsection{Tabla de síntesis}

\begin{table}[H]
\centering
\begin{tabular}{lp{8cm}p{4cm}}
\toprule
Sección & Enfoque & Recomendación posterior \\
\midrule
Repaso de los fundamentos del MDP & Relación entre la iteración de Bellman y la iteración de política/value & Revisar de nuevo el teorema de contracción y sus formas de evolución \\
Q-Learning tabular & Convergencia off-policy y equilibrio de la exploración & Revisar la cobertura del replay buffer en las tareas actuales \\
Fitted Q-Iteration y DQN & Aproximación de funciones, técnicas de estabilización y detalles de optimización & Verificar en sistemas existentes el efecto de target network / double DQN \\
De DQN a Actor-Critic & Transición hacia continuous control y la arquitectura actor-critic & Comparar la estimación del gradiente de deterministic policy gradient con la de Q-learning \\
Práctica de sistemas y monitoreo & Métricas, estado del replay buffer, mecanismos de intervención & Establecer un mecanismo de instrumentation revoke \\
Variantes y costos & Rainbow, modelo de costos, pila de inferencia & Planificar el presupuesto de entrenamiento/inferencia y conservar una ruta de reincorporación \\
Casos de aplicación y gobernanza & Avances en benchmark y sim-to-real, aspectos de gobernanza & Planificar soft/hard interrupt y la reversión ante anomaly \\
\bottomrule
\end{tabular}
\caption{Contenidos principales de esta clase y puntos de entrada para la ingeniería}
\end{table}

\subsection{Lecturas para profundizar}

\begin{knowledgebox}{Material de ampliación}
\begin{itemize}
    \item Watkins \& Dayan, ``Q-Learning,'' Machine Learning 1992
    \item Mnih et al., ``Human-level Control through Deep Reinforcement Learning (DQN),'' Nature 2015
    \item Van Hasselt et al., ``Deep Reinforcement Learning with Double Q-Learning,'' AAAI 2016
    \item Lillicrap et al., ``Continuous Control with Deep Reinforcement Learning (DDPG),'' ICLR 2016
    \item Sutton \& Barto, ``Reinforcement Learning: An Introduction,'' 2nd Edition, 2018
    \item Documentación de OpenAI Safety Gym 2.0 (permite hacer pruebas de interruptibility)
\end{itemize}
\end{knowledgebox}

\subsection{Resumen de la sección}

Esta clase parte del MDP y la ecuación de Bellman para desarrollar, paso a paso, la evolución de Q-Learning: de la forma tabular a las redes profundas y, después, a Actor-Critic. A lo largo del recorrido intercala la práctica de sistemas, las variantes y su COST, así como casos de gobernanza y de aplicación, y ofrece un checklist completo de repaso y de operación.

\end{document}
